% Фрагмент тексту за темою магістерської роботи, на якому показано роботу стилів цитування. Його не підключають через \input: main.tex передає файл у csl-demo.lua, а той замінює кожну команду dcite з ключем на цитату в потрібному стилі. Ключі — ті, що Better BibTeX присвоює записам бібліотеки Zotero у zotero.bib.
\textbf{Тема:} Удосконалення методу синтезу мовлення з керуванням голосом за текстовим описом для української мови

Сучасні нейромережеві системи синтезу мовлення здатні клонувати голос диктора за кількома секундами референсного запису, однак сприйняття природності та схожості такого мовлення суттєво різниться між системами, зокрема ElevenLabs, StyleTTS2 та XTTS-v2 \dcite{bakkoucheWhatDeterminesSuccess2026}. Водночас висока реалістичність синтезованого мовлення створює ризики для систем голосової автентифікації, оскільки детектори синтетичного мовлення помітно втрачають точність на голосах, згенерованих комерційними системами \dcite{dzhaliukEvaluatingSpeakerVerification2026}.

Альтернативою клонуванню є керування голосом за текстовим описом, коли бажані характеристики голосу, такі як стать, вік, тембр, темп чи емоційне забарвлення, задаються природною мовою без референсного запису. Проте навіть у найкращих системах цього класу зберігається розрив між інструкцією користувача та тим, як згенероване мовлення сприймають слухачі \dcite{linYouHearWhat2026}. Для навчання таких моделей потрібні корпуси, у яких кожному аудіозапису відповідає текстовий опис голосу, і одним зі способів їх отримати є автоматичне перетворення паралінгвістичних акустичних ознак на текстові промпти \dcite{kimParalinguisticAcousticFeature2027}.

Для української мови ця задача ускладнюється обмеженістю відповідних даних. Багатомовні моделі під час міжмовного клонування схильні переносити мовні та діалектні ознаки референсного запису на цільову мову \dcite{dorjeeDiagnosingReferenceSideLanguage2026}, а в умовах малого обсягу та низької якості навчальних даних zero-shot синтез на основі архітектур VITS і StyleTTS2 розглядають як засіб доповнення навчальної вибірки \dcite{hanzlicekLeveragingZeroShotTTS2027}. Тому метою роботи є удосконалення методу синтезу мовлення з керуванням голосом за текстовим описом шляхом його адаптації до української мови.
